% =============================================================================
%  SPECTRA — Section III: Problem Formulation
%  Target: IEEE Transactions on Information Forensics & Security (TIFS)
%  Author: Sagar Korde
% =============================================================================
%
%  Required packages (add to preamble if not already present):
%    \usepackage{amsmath, amssymb, amsthm}
%    \usepackage{mathtools}
%    \usepackage{bm}
%    \usepackage{booktabs}
%    \usepackage{siunitx}
%    \usepackage{thmtools}
%    \usepackage{mdframed}   % for the boxed problem statement
%
%  Theorem-like environments (add to preamble):
%    \theoremstyle{definition}
%    \newtheorem{definition}{Definition}[section]
%    \newtheorem{problem}[definition]{Problem}
% =============================================================================

\section{Problem Formulation}
\label{sec:problem}

We formalize the \textsc{SPECTRA} pipeline as a sequence of well-defined
mathematical objects.  Table~\ref{tab:notation} summarizes all principal
symbols used throughout the remainder of the paper.

% ---------------------------------------------------------------------------
%  Notation table
% ---------------------------------------------------------------------------
\begin{table}[!t]
  \renewcommand{\arraystretch}{1.18}
  \caption{Principal Notation}
  \label{tab:notation}
  \centering
  \small
  \begin{tabular}{@{}lll@{}}
    \toprule
    \textbf{Symbol} & \textbf{Description} & \textbf{Value / Dimension} \\
    \midrule
    $\mathcal{G}$ & Bitcoin address graph & --- \\
    $\mathcal{V}$ & Raw wallet-address node set & $|\mathcal{V}|=\num{2879144}$ \\
    $\mathcal{V}^{*}$ & Hub-pruned node set & $|\mathcal{V}^{*}|=\num{30000}$ \\
    $\mathcal{E}$ & Directed edge set (sender$\to$receiver) & --- \\
    $w_{uv}$ & Edge weight (BTC transferred) & $\mathbb{R}_{\geq 0}$ \\
    $N$ & Observation window (transactions) & $\num{300000}$ \\
    $\mathbf{A}_{\mathrm{sym}}$ & Symmetrized adjacency matrix & $\mathbb{R}^{|\mathcal{V}^{*}|\times|\mathcal{V}^{*}|}$ \\
    $\mathbf{L}_{\mathrm{sym}}$ & Normalized graph Laplacian & $\mathbb{R}^{|\mathcal{V}^{*}|\times|\mathcal{V}^{*}|}$ \\
    $\boldsymbol{\Phi}$ & Spectral embedding matrix & $\mathbb{R}^{|\mathcal{V}^{*}|\times k},\;k=10$ \\
    $\mathbf{f}_{v}$ & Hand-crafted node features & $\mathbb{R}^{6}$ \\
    $\mathbf{X}_{\mathrm{node}}$ & Augmented node feature matrix & $\mathbb{R}^{\num{30000}\times 16}$ \\
    $\mathbf{x}$ & Tabular transaction feature vector & $\mathbb{R}^{20}$ \\
    $\mathbf{z}$ & VAE latent code & $\mathbb{R}^{32}$ \\
    $\mathbf{Z}_{\mathrm{all}}$ & Full latent matrix & $\mathbb{R}^{N\times 32}$ \\
    $r(\mathbf{x})$ & VAE reconstruction error (anomaly score) & $\mathbb{R}_{\geq 0}$ \\
    $\tau_{\mathrm{rec}}$ & Reconstruction anomaly threshold & $0.1694$ \\
    $\mathbf{P}_{\mathrm{soft}}$ & HDBSCAN soft membership & $\mathbb{R}^{N\times 135}$ \\
    $\tau_{v}$ & Bayesian trust score for address $v$ & $[0,1]$ \\
    $\gamma_{v}$ & Markov anomaly score for address $v$ & $\mathbb{R}_{\geq 0}$ \\
    $\hat{y}_{i}$ & Predicted transaction-type label & $\{0,\ldots,5\}$ \\
    $\mathcal{A}_{t}$ & Drift alert flag at window $t$ & $\{0,1\}$ \\
    $\delta$ & Wasserstein alert threshold & $0.1$ \\
    \bottomrule
  \end{tabular}
\end{table}

% ---------------------------------------------------------------------------
%  III-A  Bitcoin Transaction Graph
% ---------------------------------------------------------------------------
\subsection{Bitcoin Transaction Graph}
\label{subsec:graph}

\begin{definition}[Bitcoin Address Graph]
\label{def:graph}
Let $\mathcal{G} = (\mathcal{V},\,\mathcal{E},\,\mathbf{W})$ be a
\emph{directed weighted graph} in which:
\begin{itemize}
  \item $\mathcal{V}$ is the set of unique Bitcoin wallet addresses observed
        over an observation window of $N = \num{300000}$ transactions,
        $|\mathcal{V}| = \num{2879144}$;
  \item $\mathcal{E} \subseteq \mathcal{V}\times\mathcal{V}$ is the set of
        directed edges, where $(u,v)\in\mathcal{E}$ if address $u$ appears as
        a sender and address $v$ appears as a receiver in at least one
        transaction; and
  \item $\mathbf{W}:\mathcal{E}\to\mathbb{R}_{>0}$ is the weight function,
        with $w_{uv}$ equal to the total BTC transferred from $u$ to $v$
        summed across all transactions in the window.
\end{itemize}
\end{definition}

\noindent\textbf{Address parsing.}
Each raw transaction record exposes two semicolon-delimited
\texttt{VARCHAR[]} arrays: \texttt{input\_addresses} (senders) and
\texttt{output\_addresses} (receivers).
For a transaction $t_i$ with value $\mathrm{val}_i$, sender set
$\mathcal{I}_i$, and receiver set $\mathcal{O}_i$, each directed edge
$(u,v)\in\mathcal{I}_i\times\mathcal{O}_i$ receives an additive weight
contribution of
\begin{equation}
  \Delta w_{uv}^{(i)}
    \;=\;
    \frac{\mathrm{val}_i}{|\mathcal{I}_i|\;\cdot\;|\mathcal{O}_i|},
  \label{eq:edge_weight}
\end{equation}
so that the total BTC value of every transaction is exactly preserved in the
edge-weight sum regardless of the number of inputs or outputs.

\noindent\textbf{Hub-node pruning.}
The raw graph $\mathcal{G}$ contains $|\mathcal{V}| = \num{2879144}$ nodes,
making full spectral decomposition computationally intractable.
We therefore retain only the $|\mathcal{V}^{*}| = \num{30000}$ nodes with
the highest \emph{undirected} degree in $\mathcal{G}$, forming the pruned
hub graph $\mathcal{G}^{*} = (\mathcal{V}^{*}, \mathcal{E}^{*}, \mathbf{W}^{*})$
where $\mathcal{E}^{*}$ and $\mathbf{W}^{*}$ are the restrictions of
$\mathcal{E}$ and $\mathbf{W}$ to $\mathcal{V}^{*}\times\mathcal{V}^{*}$.
Hub nodes collectively account for the dominant fraction of transaction volume
and preserve the structural skeleton of the network~\cite{leskovec2007graph}.

% ---------------------------------------------------------------------------
%  III-B  Spectral Embedding
% ---------------------------------------------------------------------------
\subsection{Spectral Graph Embedding}
\label{subsec:spectral}

\begin{definition}[Normalized Graph Laplacian]
\label{def:laplacian}
Given the pruned graph $\mathcal{G}^{*}$, let
$\mathbf{A}\in\mathbb{R}^{|\mathcal{V}^{*}|\times|\mathcal{V}^{*}|}$ be its
weighted adjacency matrix.  Define:
\begin{align}
  \mathbf{A}_{\mathrm{sym}}
    &= \frac{\mathbf{A} + \mathbf{A}^{\!\top}}{2},
    \label{eq:Asym} \\[4pt]
  \mathbf{D}
    &= \operatorname{diag}\!\bigl(\mathbf{A}_{\mathrm{sym}}\,\mathbf{1}\bigr),
    \label{eq:degree} \\[4pt]
  \mathbf{L}
    &= \mathbf{D} - \mathbf{A}_{\mathrm{sym}},
    \label{eq:laplacian} \\[4pt]
  \mathbf{L}_{\mathrm{sym}}
    &= \mathbf{D}^{-1/2}\,\mathbf{L}\,\mathbf{D}^{-1/2},
    \label{eq:Lsym}
\end{align}
where $\mathbf{1}$ denotes the all-ones vector.
$\mathbf{L}_{\mathrm{sym}}$ is the \emph{normalized (symmetric)} graph
Laplacian with eigenvalues in $[0,2]$.
\end{definition}

\noindent\textbf{Spectral embedding.}
Let $0 = \lambda_1 \leq \lambda_2 \leq \cdots \leq \lambda_{|\mathcal{V}^{*}|}$
be the eigenvalues of $\mathbf{L}_{\mathrm{sym}}$ and
$\mathbf{u}_1, \mathbf{u}_2, \ldots$ the corresponding orthonormal
eigenvectors.
The spectral embedding is the matrix formed by the bottom-$k$
\emph{non-trivial} eigenvectors:
\begin{equation}
  \boldsymbol{\Phi}
    = \bigl[\mathbf{u}_{2},\,\mathbf{u}_{3},\,\ldots,\,\mathbf{u}_{k+1}\bigr]
    \;\in\; \mathbb{R}^{|\mathcal{V}^{*}|\times k},
  \quad k = 10.
  \label{eq:Phi}
\end{equation}
Each row $\boldsymbol{\phi}_{v}\in\mathbb{R}^{k}$ encodes the spectral
position of node $v$ in the graph.
Randomized truncated SVD is used to compute~\eqref{eq:Phi} efficiently.

\noindent\textbf{Augmented node features.}
For each node $v\in\mathcal{V}^{*}$ we compute a six-dimensional hand-crafted
feature vector
\begin{equation}
  \mathbf{f}_{v}
    = \bigl[
        f^{\mathrm{fi}}_v,\;
        f^{\mathrm{vel}}_v,\;
        f^{\mathrm{arv}}_v,\;
        f^{\mathrm{asv}}_v,\;
        f^{\mathrm{tgm}}_v,\;
        f^{\mathrm{tgs}}_v
      \bigr]^{\!\top} \in \mathbb{R}^{6},
  \label{eq:fv}
\end{equation}
where the components represent the \emph{fan-in ratio}
($f^{\mathrm{fi}}$), \emph{transaction velocity}
($f^{\mathrm{vel}}$), \emph{average received value}
($f^{\mathrm{arv}}$), \emph{average sent value}
($f^{\mathrm{asv}}$), \emph{mean temporal gap}
($f^{\mathrm{tgm}}$), and \emph{standard deviation of temporal gap}
($f^{\mathrm{tgs}}$), respectively.
The augmented feature vector for node $v$ is the concatenation
\begin{equation}
  \mathbf{h}_{v}
    = \bigl[\boldsymbol{\phi}_{v};\; \mathbf{f}_{v}\bigr]
    \;\in\; \mathbb{R}^{16},
  \label{eq:hv}
\end{equation}
yielding the full node feature matrix
$\mathbf{X}_{\mathrm{node}}\in\mathbb{R}^{\num{30000}\times 16}$.

% ---------------------------------------------------------------------------
%  III-C  Variational Autoencoder Latent Space
% ---------------------------------------------------------------------------
\subsection{Variational Autoencoder Latent Space}
\label{subsec:vae}

\begin{definition}[$\beta$-VAE Anomaly Model]
\label{def:vae}
Let $\mathbf{x}\in\mathbb{R}^{20}$ be a min-max scaled tabular transaction
feature vector.  A $\beta$-Variational Autoencoder comprises:
\begin{enumerate}
  \item An \emph{encoder} $q_{\phi}(\mathbf{z}|\mathbf{x}) =
        \mathcal{N}\!\bigl(\boldsymbol{\mu}_{\phi}(\mathbf{x}),\;
        \operatorname{diag}(\boldsymbol{\sigma}^{2}_{\phi}(\mathbf{x}))\bigr)$
        with latent dimension $d=32$; and
  \item A \emph{decoder} $p_{\theta}(\mathbf{x}|\mathbf{z})$ that
        reconstructs $\hat{\mathbf{x}}$ from a sampled code
        $\mathbf{z}\in\mathbb{R}^{32}$.
\end{enumerate}
\end{definition}

\noindent\textbf{Training objective.}
Parameters $\phi,\theta$ are jointly learned by maximizing the
Evidence Lower BOund (ELBO):
\begin{equation}
  \mathcal{L}_{\mathrm{ELBO}}(\phi,\theta;\mathbf{x})
    = \mathbb{E}_{q_{\phi}}\!\bigl[\log p_{\theta}(\mathbf{x}|\mathbf{z})\bigr]
      - \beta\,\mathrm{KL}\!\Bigl(q_{\phi}(\mathbf{z}|\mathbf{x})
        \,\big\|\,p(\mathbf{z})\Bigr),
  \label{eq:elbo}
\end{equation}
where $p(\mathbf{z})=\mathcal{N}(\mathbf{0},\mathbf{I})$ and $\beta = 1.0$.

\noindent\textbf{Anomaly scoring.}
After training, the per-sample reconstruction error
\begin{equation}
  r(\mathbf{x})
    = \bigl\|\mathbf{x} - \hat{\mathbf{x}}\bigr\|_{2}^{2}
  \label{eq:recon}
\end{equation}
serves as a continuous anomaly score.
A transaction is flagged as anomalous if
$r(\mathbf{x}) > \tau_{\mathrm{rec}}$, where the threshold
\begin{equation}
  \tau_{\mathrm{rec}}
    = \mathrm{Percentile}_{95}\!\bigl(r(\mathbf{x})\bigr)_{\mathbf{x}\in\mathcal{X}_{\mathrm{train}}}
    = 0.1694
  \label{eq:tau_rec}
\end{equation}
is set to the 95th percentile of reconstruction errors on the training set,
ensuring a nominal false-positive rate of $5\%$.
The full latent matrix
$\mathbf{Z}_{\mathrm{all}} = \{\mathbf{z}_{i}\}_{i=1}^{N}
\in\mathbb{R}^{N\times 32}$
is retained for downstream density clustering.

% ---------------------------------------------------------------------------
%  III-D  HDBSCAN Density Clustering
% ---------------------------------------------------------------------------
\subsection{HDBSCAN Density Clustering}
\label{subsec:hdbscan}

\begin{definition}[Latent-Space Cluster Assignment]
\label{def:hdbscan}
Given $\mathbf{Z}_{\mathrm{all}}\in\mathbb{R}^{N\times 32}$, hierarchical
density-based spatial clustering of applications with noise
(HDBSCAN~\cite{campello2013density}) assigns each transaction a label
\begin{equation}
  c_{i} \;\in\; \{-1,\,0,\,1,\,\ldots,\,C-1\},
  \label{eq:cluster_label}
\end{equation}
where $c_{i}=-1$ denotes a noise point and $C$ denotes the number of
discovered clusters.
\end{definition}

\noindent\textbf{Empirical result.}
Applied to our dataset, HDBSCAN discovers $C = 135$ clusters with
\num{63390} noise points ($21.1\%$ of all transactions).
HDBSCAN additionally provides \emph{soft membership probabilities}
\begin{equation}
  \mathbf{P}_{\mathrm{soft}}
    \in [0,1]^{N\times C},
  \qquad
  \bigl(\mathbf{P}_{\mathrm{soft}}\bigr)_{ic}
    = \Pr[c_{i} = c \mid \mathbf{z}_{i}],
  \label{eq:soft}
\end{equation}
which are propagated to downstream modules as probabilistic cluster
evidence rather than hard assignments.

% ---------------------------------------------------------------------------
%  III-E  Wasserstein Drift Detection
% ---------------------------------------------------------------------------
\subsection{Wasserstein Drift Detection}
\label{subsec:wasserstein}

\begin{definition}[Temporal Drift Score]
\label{def:drift}
Partition the $N$ transactions into $T=30$ non-overlapping monthly time
windows $\{\mathcal{W}_{t}\}_{t=1}^{T}$.
For each window $t$, let $D_t$ be the empirical distribution of cluster
labels $\{c_{i}\}_{i\in\mathcal{W}_{t}}$.
The \emph{Wasserstein drift score} between consecutive windows is
\begin{equation}
  d_{t}
    = \mathcal{W}_{1}(D_{t-1},\, D_{t}),
  \quad t = 2,\ldots,T,
  \label{eq:wasserstein}
\end{equation}
where $\mathcal{W}_{1}(\cdot,\cdot)$ denotes the Wasserstein-$1$
(earth-mover's) distance.  A drift alert is raised whenever
$d_{t} > \delta$, i.e.,
\begin{equation}
  \mathcal{A}_{t} = \mathbf{1}\!\bigl[d_{t} > \delta\bigr],
  \quad \delta = 0.1.
  \label{eq:alert}
\end{equation}
\end{definition}

\noindent\textbf{Empirical result.}
All $T-1 = 29$ consecutive window pairs yield $\mathcal{A}_{t}=1$, with
drift scores ranging from $0.179$ to $50.745$ (mean $8.500$), indicating
persistent and substantial distributional shift in the Bitcoin transaction
stream throughout the observation period.

% ---------------------------------------------------------------------------
%  III-F  Bayesian Trust Scoring
% ---------------------------------------------------------------------------
\subsection{Bayesian Trust Scoring}
\label{subsec:trust}

\begin{definition}[Bayesian Trust Score]
\label{def:trust}
For each wallet address $v$, let
$\mathcal{T}_{v} = \{\mathbf{x}_{i}\}$ be the set of transactions
associated with $v$.
We model the \emph{trustworthiness} of $v$ as a latent probability
$\rho_{v}\sim\mathrm{Beta}(\alpha,\beta)$ and update a Beta posterior
sequentially over the transactions attributed to $v$.
\end{definition}

\noindent\textbf{Prior.}
We adopt a symmetric, weakly informative prior:
$\alpha_{0} = \beta_{0} = 2$, encoding mild belief that addresses are
neither universally trustworthy nor universally anomalous.

\noindent\textbf{Likelihood.}
For each transaction $i\in\mathcal{T}_{v}$, define the \emph{likelihood
weight}
\begin{equation}
  \ell_{i}
    = \exp\!\bigl(-\kappa\,r(\mathbf{x}_{i})\bigr),
  \quad \kappa = 1.0,
  \label{eq:likelihood}
\end{equation}
where $r(\mathbf{x}_{i})$ is the VAE reconstruction error from
\eqref{eq:recon}.  Transactions with low anomaly score contribute near-unity
weight, while high-anomaly transactions contribute near-zero weight.

\noindent\textbf{Posterior update.}
The Beta posterior parameters after observing all transactions in
$\mathcal{T}_{v}$ are
\begin{align}
  \alpha_{n} &= \alpha_{0} + \textstyle\sum_{i\in\mathcal{T}_{v}} \ell_{i},
  \label{eq:alpha_n} \\
  \beta_{n}  &= \beta_{0}  + \textstyle\sum_{i\in\mathcal{T}_{v}} (1 - \ell_{i}).
  \label{eq:beta_n}
\end{align}
The \emph{trust score} of address $v$ is the posterior mean:
\begin{equation}
  \tau_{v}
    = \mathbb{E}\!\bigl[\mathrm{Beta}(\alpha_{n},\beta_{n})\bigr]
    = \frac{\alpha_{n}}{\alpha_{n} + \beta_{n}}
    \;\in\; (0,1).
  \label{eq:trust}
\end{equation}
Across $|\mathcal{V}^{*}|=\num{30000}$ hub addresses, the empirical trust
distribution has mean $0.634$ and standard deviation $0.105$.

% ---------------------------------------------------------------------------
%  III-G  Markov Chain Behavioral Profiling
% ---------------------------------------------------------------------------
\subsection{Markov Chain Behavioral Profiling}
\label{subsec:markov}

\begin{definition}[Address Behavioral State]
\label{def:markov}
Associate with each transaction $i$ a \emph{behavioral state}
$s_{i}\in\mathcal{S}$, where
\begin{equation}
  \mathcal{S}
    = \{0,\,1,\,\ldots,\,C-1\} \cup \{\text{noise}\},
  \quad |\mathcal{S}| = 136,
  \label{eq:statespace}
\end{equation}
obtained from the HDBSCAN cluster label $c_{i}$ (with $c_{i}=-1$ mapped
to the designated noise state).
For each address $v$, the sequence of states
$s_{(1)}^{v}, s_{(2)}^{v}, \ldots$ in chronological order defines a
\emph{behavioral trajectory}.
\end{definition}

\noindent\textbf{Transition matrix.}
The global transition probability matrix
$\mathbf{P}\in\mathbb{R}^{|\mathcal{S}|\times|\mathcal{S}|}$
is estimated from all observed consecutive-state pairs with Laplace smoothing:
\begin{equation}
  P_{jk}
    = \frac{n_{jk} + \varepsilon}{\displaystyle\sum_{k'} (n_{jk'} + \varepsilon)},
  \quad \varepsilon = 10^{-6},
  \label{eq:transition}
\end{equation}
where $n_{jk}$ counts all transitions from state $j$ to state $k$ across all
addresses.

\noindent\textbf{Stationary distribution.}
The stationary distribution $\boldsymbol{\pi}\in\mathbb{R}^{|\mathcal{S}|}$
satisfies $\mathbf{P}^{\!\top}\boldsymbol{\pi} = \boldsymbol{\pi}$ and is
computed via power iteration with tolerance $10^{-8}$.

\noindent\textbf{Anomaly score.}
For each address $v$, let $(s_{t-1}^{v}, s_{t}^{v})$ denote the most recently
observed state transition.  The \emph{Markov anomaly score} is defined as
\begin{equation}
  \gamma_{v}
    = -\log P_{s_{t-1}^{v},\,s_{t}^{v}},
  \label{eq:markov_score}
\end{equation}
so that rare transitions yield large values.
Across hub addresses, the empirical distribution of $\gamma_{v}$ has mean
$0.875$ and maximum $1.058$.

% ---------------------------------------------------------------------------
%  III-H  GraphSAGE Classification
% ---------------------------------------------------------------------------
\subsection{GraphSAGE Node Classification}
\label{subsec:graphsage}

\begin{definition}[Composite Node Feature]
\label{def:node_feature}
For each hub node $v\in\mathcal{V}^{*}$, define the composite input feature
vector
\begin{equation}
  \tilde{\mathbf{h}}_{v}^{(0)}
    = \bigl[
        \mathbf{h}_{v};\;
        \tau_{v};\;
        \gamma_{v};\;
        \mathbf{e}_{c_{v}}
      \bigr]
    \;\in\; \mathbb{R}^{d_{\mathrm{in}}},
  \label{eq:composite}
\end{equation}
where $\mathbf{h}_{v}\in\mathbb{R}^{16}$ is the spectral-augmented feature
from~\eqref{eq:hv}, $\tau_{v}\in(0,1)$ is the Bayesian trust score
from~\eqref{eq:trust}, $\gamma_{v}\in\mathbb{R}_{\geq 0}$ is the Markov
anomaly score from~\eqref{eq:markov_score}, and
$\mathbf{e}_{c_{v}}\in\{0,1\}^{\min(C,20)}$ is a one-hot cluster indicator
(truncated to the 20 most frequent clusters for dimensionality control).
Thus $d_{\mathrm{in}} = 16 + 1 + 1 + \min(135, 20) = 38$.
\end{definition}

\noindent\textbf{GraphSAGE propagation.}
A two-layer GraphSAGE~\cite{hamilton2017inductive} network iteratively
aggregates neighbor information:
\begin{equation}
  \mathbf{h}_{v}^{(l+1)}
    = \sigma\!\Biggl(
        \mathbf{W}^{(l)}\!
        \cdot
        \mathrm{MEAN}\!\Bigl(
          \bigl\{\mathbf{h}_{v}^{(l)}\bigr\}
          \cup
          \bigl\{\mathbf{h}_{u}^{(l)} : u\in\mathcal{N}(v)\bigr\}
        \Bigr)
      \Biggr),
  \label{eq:graphsage}
\end{equation}
where $\mathcal{N}(v)$ denotes the neighborhood of $v$ in $\mathcal{G}^{*}$,
$\mathbf{W}^{(l)}$ is a learnable weight matrix, and $\sigma$ is the ReLU
activation.  The final node embedding $\mathbf{h}_{v}^{(2)}$ is passed
through a softmax head to predict the transaction-type label
$\hat{y}_{v}\in\{0,1,\ldots,5\}$ over six classes.

\noindent\textbf{Training.}
Parameters are optimized with cross-entropy loss and the Adam optimizer,
with early stopping (patience $= 15$ epochs) applied on a held-out
validation split to prevent overfitting.

% ---------------------------------------------------------------------------
%  III-I  Formal Problem Statement
% ---------------------------------------------------------------------------
\subsection{Formal Problem Statement}
\label{subsec:problem}

We are now in a position to state the central problem addressed by SPECTRA.

\begin{mdframed}[linewidth=0.8pt, innertopmargin=8pt, innerbottommargin=8pt,
                 innerleftmargin=10pt, innerrightmargin=10pt]
\begin{problem}[SPECTRA Inference Problem]
\label{prob:spectra}
\textbf{Given:}
\begin{enumerate}
  \item A directed weighted Bitcoin address graph
        $\mathcal{G}=(\mathcal{V},\mathcal{E},\mathbf{W})$
        constructed over $N=\num{300000}$ transactions, pruned to
        hub subgraph $\mathcal{G}^{*}$ with $|\mathcal{V}^{*}|=\num{30000}$
        nodes;
  \item A tabular transaction feature matrix
        $\mathbf{X}\in\mathbb{R}^{N\times 20}$; and
  \item A temporal sequence $\mathcal{T} = (t_{1}, t_{2}, \ldots, t_{N})$
        of transaction timestamps partitioned into $T=30$ monthly windows.
\end{enumerate}

\medskip
\textbf{Find:}
\begin{enumerate}
  \item[\textbf{(F1)}] A cluster assignment
        $c_{i}\in\{-1,0,\ldots,C-1\}$ for each transaction $i$, together
        with soft membership probabilities $\mathbf{P}_{\mathrm{soft}}$,
        that reflect latent behavioral patterns in the VAE latent space;
  \item[\textbf{(F2)}] A trust score $\tau_{v}\in(0,1)$ for each hub address
        $v\in\mathcal{V}^{*}$, certifying the aggregate anomaly level of
        all transactions attributed to $v$;
  \item[\textbf{(F3)}] A transaction-type prediction
        $\hat{y}_{i}\in\{0,\ldots,5\}$ for each transaction $i$, produced
        by the GraphSAGE classifier trained on the composed node features;
        and
  \item[\textbf{(F4)}] A sequence of binary drift alert flags
        $\{\mathcal{A}_{t}\}_{t=2}^{T}$, signaling statistically significant
        concept drift in the cluster-label distribution between consecutive
        monthly windows.
\end{enumerate}

\medskip
\textbf{Subject to:}
\begin{enumerate}
  \item[\textbf{(S1)}] \emph{Temporal fairness.} No temporal identifiers
        (timestamps, block heights, or chronological indices) are used as
        input features to \textbf{(F3)}, preventing the model from learning
        spurious correlations with annotator conventions tied to particular
        time periods.
  \item[\textbf{(S2)}] \emph{Inductive generalization.} The GraphSAGE
        classifier \textbf{(F3)} and the trust-scoring module \textbf{(F2)}
        must generalize to wallet addresses unseen during training, relying
        solely on the structural and latent-space features derivable from
        the neighborhood of a new node in $\mathcal{G}$.
\end{enumerate}
\end{problem}
\end{mdframed}

\noindent
Together, \textbf{(F1)}--\textbf{(F4)} constitute a \emph{multi-objective}
inference problem whose outputs are jointly coherent: the cluster assignments
inform the node features, the trust scores encode anomaly evidence, and the
drift alerts provide actionable runtime signals about distribution shift.
The next section details the SPECTRA architecture that realizes this
formulation.
